\documentclass[10pt,a4paper]{amsart}
\usepackage{amsmath,amssymb,amsthm,graphicx}
\usepackage{verbatim}
\newenvironment{reference}{\comment}{\endcomment}
%\newenvironment{reference}{}{}
\newenvironment{slogan}{\comment}{\endcomment}
\newenvironment{history}{\comment}{\endcomment}

% For commutative diagrams we use Xy-pic
\usepackage[all]{xy}

% We use 2cell for 2-commutative diagrams.
\xyoption{2cell}
% Theorem environments.
%
\theoremstyle{plain}
\newtheorem{theorem}[subsection]{Theorem}
\newtheorem{proposition}[subsection]{Proposition}
\newtheorem{lemma}[subsection]{Lemma}

\theoremstyle{definition}
\newtheorem{definition}[subsection]{Definition}
\newtheorem{example}[subsection]{Example}
\newtheorem{exercise}[subsection]{Exercise}
\newtheorem{situation}[subsection]{Situation}

\theoremstyle{remark}
\newtheorem{remark}[subsection]{Remark}
\newtheorem{remarks}[subsection]{Remarks}

\numberwithin{equation}{subsection}

% Macros
%
\def\lim{\mathop{\mathrm{lim}}\nolimits}
\def\colim{\mathop{\mathrm{colim}}\nolimits}
\def\Spec{\mathop{\mathrm{Spec}}}
\def\Hom{\mathop{\mathrm{Hom}}\nolimits}
\def\Ext{\mathop{\mathrm{Ext}}\nolimits}
\def\SheafHom{\mathop{\mathcal{H}\!\mathit{om}}\nolimits}
\def\SheafExt{\mathop{\mathcal{E}\!\mathit{xt}}\nolimits}
\def\Sch{\mathit{Sch}}
\def\Mor{\mathop{\mathrm{Mor}}\nolimits}
\def\Ob{\mathop{\mathrm{Ob}}\nolimits}
\def\Sh{\mathop{\mathit{Sh}}\nolimits}
\def\NL{\mathop{N\!L}\nolimits}
\def\CH{\mathop{\mathrm{CH}}\nolimits}
\def\proetale{{pro\text{-}\acute{e}tale}}
\def\etale{{\acute{e}tale}}
\def\QCoh{\mathit{QCoh}}
\def\Ker{\mathop{\mathrm{Ker}}}
\def\Im{\mathop{\mathrm{Im}}}
\def\Coker{\mathop{\mathrm{Coker}}}
\def\Coim{\mathop{\mathrm{Coim}}}

% Boxtimes
%
\DeclareMathSymbol{\boxtimes}{\mathbin}{AMSa}{"02}

%
% Macros for moduli stacks/spaces
%
\def\QCohstack{\mathcal{QC}\!\mathit{oh}}
\def\Cohstack{\mathcal{C}\!\mathit{oh}}
\def\Spacesstack{\mathcal{S}\!\mathit{paces}}
\def\Quotfunctor{\mathrm{Quot}}
\def\Hilbfunctor{\mathrm{Hilb}}
\def\Curvesstack{\mathcal{C}\!\mathit{urves}}
\def\Polarizedstack{\mathcal{P}\!\mathit{olarized}}
\def\Complexesstack{\mathcal{C}\!\mathit{omplexes}}
% \Pic is the operator that assigns to X its picard group, usage \Pic(X)
% \Picardstack_{X/B} denotes the Picard stack of X over B
% \Picardfunctor_{X/B} denotes the Picard functor of X over B
\def\Pic{\mathop{\mathrm{Pic}}\nolimits}
\def\Picardstack{\mathcal{P}\!\mathit{ic}}
\def\Picardfunctor{\mathrm{Pic}}
\def\Deformationcategory{\mathcal{D}\!\mathit{ef}}
\usepackage{cleveref}
\crefname{originalproposition}{original proposition}{original propositions}
\Crefname{originalproposition}{Original proposition}{Original propositions}
\makeatletter
\newlabel{proposition-compute-bivariant}{{0FBT}{source}{}{sourcepointer-1}{}}
\newlabel{proposition-compute-bivariant@cref}{{[originalproposition][0][]0FBT}{[1][0][]source}{}{}{}}
\makeatother
\crefname{originallemma}{original lemma}{original lemmas}
\Crefname{originallemma}{Original lemma}{Original lemmas}
\makeatletter
\newlabel{lemma-chow-cohomology-towards-point}{{0FBW}{source}{}{sourcepointer-2}{}}
\newlabel{lemma-chow-cohomology-towards-point@cref}{{[originallemma][0][]0FBW}{[1][0][]source}{}{}{}}
\makeatother
\title{For a smooth equidimensional scheme over a field the bivariant Chow ring is Chow homology and composition becomes intersection}
\author{The Stacks Project Authors}
\date{stacks-project:chow.tex:local-snapshot-20261012; GNU FDL 1.2 or later}
\begin{document}
\maketitle
\noindent\textit{Editorial note.} All delivered mathematical text is the author's own, copied from the pinned Stacks Project chapter file chow.tex; no statement or proof step was rewritten, reordered or removed.\par\medskip
\noindent\textit{Editorial note.} Preamble provenance: the pinned chapter file opens with an input of the shared preamble and carries no preamble text of its own, so the project's single Stacks-wide preamble staged at sources/stacks-chow/preamble.tex is excerpted verbatim as two spans, lines 14-18 (Xy-pic setup) and lines 154-217 (the native theorem declarations and mathematical macros). Nothing was added, removed or redefined; the wrapper class and the amsmath/amssymb/amsthm/graphicx base packages are supplied by the assembler.\par\medskip
\noindent\textit{Editorial note.} The counted statement and proof contain no cross reference, no citation and no included file, so the delivered body is self contained.\par\medskip
\noindent\textit{Editorial note.} Canonical location of the counted proposition: Stacks tag 0FC2 (sources/stacks-chow/pinned-original-source.tex line 14403).\par\medskip
\noindent\textit{Allowed standard prerequisites (editorial).}\begin{itemize}
\item The isomorphism between the bivariant ring and the Chow ring via the structure morphism. Same chapter, Proposition labelled proposition-compute-bivariant, cited at proof lines 14416-14420 and 14426-14430. Supplies the first isomorphism and the description of its inverse.
\item The identification of bivariant classes over a point with Chow groups. Same chapter, Lemma labelled lemma-chow-cohomology-towards-point, cited at proof lines 14418-14422. Supplies the second isomorphism and its inverse c\_alpha.
\item The intersection product defined via the diagonal class. Same chapter, the definition of the product preceding this lemma, invoked at proof lines 14430-14434. Identifies Delta\textasciicircum{}!(alpha x beta) with alpha . beta.
\end{itemize}\medskip
\ifdefined\hypertarget\hypertarget{sourcepointer-1}{}\fi
\noindent\textit{Source reference.} Original proposition 0FBT: sources/stacks-chow/pinned-original-source.tex, source lines 14003--14003. Same-chapter pointer into the author environment proposition at source line 14003; invoked at proof line 14421 of the counted proof. It is retained only as a navigation pointer; every load-bearing step is inside the delivered span.\par\medskip
\ifdefined\hypertarget\hypertarget{sourcepointer-2}{}\fi
\noindent\textit{Source reference.} Original lemma 0FBW: sources/stacks-chow/pinned-original-source.tex, source lines 14183--14183. Same-chapter pointer into the author environment lemma at source line 14183; invoked at proof line 14423 of the counted proof. It is retained only as a navigation pointer; every load-bearing step is inside the delivered span.\par\medskip
\begin{lemma}
\label{lemma-identify-chow-for-smooth}
Let $k$ be a field. Let $X$ be a smooth scheme over $k$, equidimensional
of dimension $d$. The map
$$
A^p(X) \longrightarrow \CH_{d - p}(X),\quad
c \longmapsto c \cap [X]_d
$$
is an isomorphism. Via this isomorphism composition of bivariant
classes turns into the intersection product defined above.
\end{lemma}

\begin{proof}
Denote $g : X \to \Spec(k)$ the structure morphism.
The map is the composition of the isomorphisms
$$
A^p(X) \to A^{p - d}(X \to \Spec(k)) \to \CH_{d - p}(X)
$$
The first is the isomorphism $c \mapsto c \circ g^*$ of
Proposition \ref{proposition-compute-bivariant}
and the second is the isomorphism $c \mapsto c \cap [\Spec(k)]$ of
Lemma \ref{lemma-chow-cohomology-towards-point}.
From the proof of Lemma \ref{lemma-chow-cohomology-towards-point}
we see that the inverse to the second arrow sends $\alpha \in \CH_{d - p}(X)$
to the bivariant class $c_\alpha$ which sends $\beta \in \CH_*(Y)$
for $Y$ locally of finite type over $k$
to $\alpha \times \beta$ in $\CH_*(X \times_k Y)$. From the proof of
Proposition \ref{proposition-compute-bivariant} we see the inverse
to the first arrow in turn sends $c_\alpha$ to the bivariant class
which sends $\beta \in \CH_*(Y)$ for $Y \to X$ locally of finite type
to $\Delta^!(\alpha \times \beta) = \alpha \cdot \beta$.
From this the final result of the lemma follows.
\end{proof}

\end{document}
